\runninghead{Note on Sources and Delivery}
\section*{Note on Sources and Delivery}

\noindent\textbf{Occasion, audience, and delivery estimate.}
This homily is written for an adult parish assembly celebrating the First Sunday of Advent, Year~B (\work{Ordo lectionum Missae} n.~2; U.S.\ Lectionary No.~2; Sunday, November~29, 2026, opening Liturgical Year~2027, Cycle~B / Weekday~I, in the Dioceses of the United States of America). The spoken body contains \textbf{1,478 spoken words} (1,474 raw source tokens), yielding an estimated speaking duration of roughly \textbf{10.6 to 11.8~minutes} at an unhurried parish pace of 125--140 words per minute (about 11.4~minutes at 130 words per minute). This figure is a word-count pace estimate, not a timed human delivery or an audible rehearsal. The homily concludes as spoken preaching without composing a new liturgical prayer.

\medskip
\noindent\textbf{Relation to the reviewed studies and decisive textual detail.}
The preacher's own plain voice carries the exegetical argument across the readings, naming only St.~John Chrysostom, St.~Gregory the Great, and St.~Augustine naturally in the spoken body while drawing its governing spine from the first reviewed interpretation of the companion studies (\emph{Divine Forbearance and the Unveiling of an Already-Present Lord}): St.~Jerome (\work{Commentarii in Isaiam} XVII, PL~24 cols.~642--651) on God's great forbearance (\emph{nimia patientia}, Isa.~63:17), the abrupt grammatical shift from wish (\emph{Utinam dirumperes}) to accomplished past tense (\emph{Descendisti}, Isa.~64:1, 3b~Vg.\ / Heb.~63:19b, 64:2b) read of the Incarnation (col.~646), inverted anger (\emph{ordo praeposterus}, 64:5~Vg.\ [Heb.~64:4]), and the lament over God's hidden face (64:7~Vg.\ [Heb.~64:6], where Jerome at col.~650 quotes Ps.~80[79]:4, \emph{Ostende faciem tuam, et salvi erimus}, ``shew us thy face, and we shall be saved,'' without quoting \emph{Deus, converte nos} or referencing a postconciliar refrain); St.~Augustine (\work{Enarrationes in Psalmos}, Ps.~79 §§1, 3, 9) on the Shepherd and vineyard (§1), Christ veiling his face in a cloud of flesh and weakness (§3), and restoring the vineyard upon the Son of man as sole foundation (§9); St.~John Chrysostom (\work{Homilies on First Corinthians}, Homily~II §6) on \emph{apokalupsis} (``unveiling,'' 1~Cor.~1:7) signifying a Lord already present though unseen; St.~Hilary of Poitiers (\work{De Trinitate} IX.67, reasoning on Mark~13:32 / Matt.~24:36 alongside Matt.~24:44, 46, applied directly to Mark~13:33 by St.~Thomas Aquinas in the \work{Catena Aurea in Marcum} 13:32--37) on why Christ withholds the hour; and St.~Bede (\work{In Marci Evangelium Expositio} IV.13, drawing on St.~Gregory the Great, \work{Homiliae in Evangelia} Homily~9 on Matt.~25:14--30) on Christ's Ascension carrying human flesh into heaven as a ``far country'' while never withdrawing his divine presence.

To keep human guilt, grace, and concrete watchfulness intact, the homily draws compatibly on non-conflicting elements from the second and third reviewed interpretations: Chrysostom (Homily~II §§7--8) on 1~Corinthians~1:8--9 warning a wavering church not to pull against God's rein; Bede (IV.13) identifying the servants' \emph{authority over every work} (Mark~13:34) with the grace of the Holy Spirit, the \emph{doorkeeper} (\emph{janitor}) with pastors, and \emph{to all} (13:37) with every believer guarding the heart; St.~Gregory the Great's three marks of vigilance (\work{Homiliae in Evangelia} Homily~13 on Luke~12:35--40, quoted by Bede IV.13); St.~Augustine (\work{Epistula}~199 §3, mediated via Aquinas's \work{Catena Aurea in Marcum} 13:32--37) explaining how the Day of the Lord meets every generation at personal death; St.~Cyril of Jerusalem (\work{Catechesis}~15 §§1, 4) on Christ's two Advents (§1) and refusing curiosity about dates (§4); St.~Augustine (\work{Enarr.}, Ps.~84 §§6, 12) on asking for ``Thy Christ'' and humbling pride (§6), and hearts furrowed by the plough of confession receiving interior \emph{sweetness} (§12); and St.~Robert Bellarmine (\work{Explanatio in Psalmos}, Ps.~84, explanation~12) on the Virgin Mary as the untouched earth yielding Christ. A genuine exegetical disagreement between St.~Jerome and St.~Thomas Aquinas documented in the companion studies is preserved here without blending: St.~Thomas Aquinas (\work{Expositio super Isaiam ad litteram}, ch.~63--64, followed by the second study interpretation) explains Isaiah~63:17 through permissive withholding of grace and prefers a historical literal sense for \emph{Descendisti} in 64:3b~Vg.\ (Heb.~64:2b; deliverance from Egypt in Exod.~4 and the cliffs of Arnon in Num.~21) rather than Jerome's direct Incarnational reading. Joining these verified verse-level patristic and Doctoral expositions across this postconciliar Year~B compilation is an editorial synthesis for preaching rather than a pre-existing sermon on this exact Mass.

The decisive textual turns changing the listener's understanding are Isaiah's leap from optative wish to accomplished past tense (Isa.~64:1, 3b~Vg.\ / Heb.~63:19b, 64:2b) set beside St.~Paul's choice of \emph{apokalupsis} (``unveiling,'' 1~Cor.~1:7) and Mark's departing householder (Mark~13:33--37). Together they overturn the assumption that God's silence means absence, disclosing an already-present Lord whose patience awaits repentance and whose nearness calls for honest confession and active watchfulness.

\medskip
\noindent\textbf{Rights, study English, numbering, and source-copy bounds.}
Approved proclaimed U.S.\ liturgical texts --- \work{The Roman Missal}, Third Typical Edition (U.S.\ English, copyright ICEL, 2010/2011), protected modern Latin Missal oration, Preface, and solemn-blessing bodies, and the U.S.\ \work{Lectionary for Mass} (New American Bible and Revised Grail / Abbey Psalms) --- are neither reproduced nor closely paraphrased. Every direct English Scripture quotation in the spoken homily comes from the public-domain Douay--Rheims (Challoner) translation strictly as \textbf{public-domain study English}, never as the approved proclaimed English of the U.S.\ Lectionary; Latin biblical controls follow the Clementine Vulgate. Hebrew and U.S.\ Lectionary Psalms~25, 80, and~85 correspond to Vulgate/Douay--Rheims Psalms~24, 79, and~84; Hebrew/Lectionary Isaiah~63:19b and~64:2b--7 correspond to Vulgate/Douay--Rheims Isaiah~64:1 and~64:3b--8.

Liturgical and copy-quality bounds verified in the companion studies govern every reference here: the \work{Gloria} is omitted on Sundays of Advent (GIRM~53); the Entrance Antiphon (\emph{Ad te levavi}, Ps.~25[24]:1--3; GIRM~48) rests on a cataloged biblical basis without offline Missal or Antiphonary truncation collation; the Communion Antiphon (\emph{Dominus dabit benignitatem}, Ps.~85[84]:13; GIRM~87) is collated against the secondary 2002 digital Missal reproduction (p.~83, with printed 2002 and printed 2008 correspondence unverified); the ancient Roman offering prayer cited is \work{Sacramentarium Veronense} XXVIIII.372 (\work{Item alia}, with no First-Sunday-of-Advent heading; Feltoe 1896, p.~74), verified against the secondary 2002 digital Missal reproduction (pp.~82--83) as matching the petition body of \emph{Suscipe, quaesumus, Domine} except for one inflection (\emph{temporalis} against \emph{temporali}), with printed 2002 and printed 2008 correspondence unverified offline and cited in the homily strictly as an ancient prayer of the Roman Church rather than a substitute for the modern Missal prayer. St.~Jerome (PL~24 cols.~642--651) and St.~Thomas Aquinas (\work{Super Isaiam}, Parma Vol.~14, pp.~370--371) were inspected in uncollated raw OCR copies; St.~Bede (\work{In Marcum} IV.13) follows the Wikisource/Corpus Corporum transcription without PL~92 print-facsimile collation; and St.~Augustine (\work{Epistula}~199 §3) is cited strictly as mediated through Aquinas's \work{Catena Aurea in Marcum}.

\section*{References}

\begin{itemize}[leftmargin=1.5em,itemsep=0.25em]

\item \textbf{Biblical Texts (Public-Domain Study Controls).} \work{Biblia Sacra juxta Vulgatam Clementinam} (eBible \texttt{latvuc}), Isa.~63:16--17, 64:1, 3--8; Ps.~24:1--3; Ps.~79:2--4, 15--16, 18--19; Ps.~84:8, 13; 1~Cor.~1:3--9; Mark~13:32--37. \work{The Holy Bible, Douay--Rheims Version} (Challoner revision, Project Gutenberg~1581), public-domain study English for the same passages (not the approved proclaimed U.S.\ Lectionary).

\item \textbf{Liturgical Books and Rubrics.} \work{Missale Romanum}, \work{editio typica tertia} (2002) and \work{reimpressio emendata} (2008), \work{Dominica I Adventus}; \work{The Roman Missal}, Third Typical Edition (U.S.\ English, 2011); protected modern Latin and ICEL English texts are not reproduced. \work{General Instruction of the Roman Missal} (U.S.\ English, 2011--2021), nn.~48, 53, 87 (summarized only). \work{Ordo lectionum Missae} (1981), n.~2; U.S.\ \work{Lectionary for Mass} (1998/2002), No.~2 (protected texts not reproduced).

\item \textbf{Ancient Roman Sacramentary.} \work{Sacramentarium Veronense} (\work{Sacramentarium Leonianum}), ed.\ C.~L.\ Feltoe (Cambridge, 1896), Section~XXVIIII, \work{Item alia}, marginal~372, p.~74 (verified bounded comparison against secondary 2002 digital Missal reproduction, pp.~82--83).

\item \textbf{St.~Augustine of Hippo.} \work{Enarrationes in Psalmos} (\work{Expositions on the Book of Psalms}), NPNF~1, Vol.~8 (1888): Ps.~79 (LXXX), §§1, 3, 9; Ps.~84 (LXXXV), §§6, 12. \work{Epistula}~199, §3, mediated via St.~Thomas Aquinas, \work{Catena Aurea in Marcum} 13:32--37.

\item \textbf{St.~Bede the Venerable.} \work{In Marci Evangelium Expositio}, Book~IV, ch.~13 (on Mark~13:32--37), Wikisource/Corpus Corporum digital Latin transcription (not print-collated against PL~92).

\item \textbf{St.~Cyril of Jerusalem.} \work{Catechesis}~15, §§1, 4, NPNF~2, Vol.~7.

\item \textbf{St.~Gregory the Great.} \work{Homiliae in Evangelia}, Homily~9 (on Matt.~25:14--30) and Homily~13 (on Luke~12:35--40), Wikisource Latin extraction.

\item \textbf{St.~Hilary of Poitiers.} \work{De Trinitate}, Book~IX, §67 (on Mark~13:32 / Matt.~24:36 alongside Matt.~24:44, 46), NPNF~2, Vol.~9 (1898).

\item \textbf{St.~Jerome.} \work{Commentarii in Isaiam Prophetam}, Book~XVII (on Isa.~63:16--17; 64:1, 3--8), Migne, \work{Patrologia Latina} Vol.~24 (Paris, 1865), cols.~642--651 (uncollated raw OCR copy).

\item \textbf{St.~John Chrysostom.} \work{Homilies on First Corinthians}, Homily~II (§§6--8, on 1~Cor.~1:7--9), NPNF~1, Vol.~12.

\item \textbf{St.~Robert Bellarmine.} \work{Explanatio in Psalmos} (trans.\ John O'Sullivan): Ps.~84, explanation~12.

\item \textbf{St.~Thomas Aquinas.} \work{Catena Aurea: Commentary on the Four Gospels}, \work{St.~Mark}, on Mark~13:32--37 (Oxford, 1842; applying Hilary's reasoning in \work{De Trinitate} IX.67 to Mark~13:33 and mediating St.~Augustine, \work{Epistula}~199 §3). \work{Expositio super Isaiam ad litteram}, ch.~63--64, Parma Vol.~14 (1863), pp.~370--371 (uncollated raw OCR copy; noted above for the contrasting literal reading of Isa.~63:17 and 64:3b~Vg.\ [Heb.~64:2b]).

\end{itemize}
